\documentclass[10pt,a4paper]{article}
\usepackage[utf8]{inputenc}
\usepackage[T1]{fontenc}
\usepackage[margin=1.8cm]{geometry}
\usepackage{enumitem}
\usepackage{booktabs}
\usepackage{parskip}
\setlength{\parskip}{0.35em}

\title{\textbf{AI-Powered Smart Agriculture Advisory System for Smallholder Farmers in Uganda}\\
\large Project Planning and Execution Readiness (February Deliverables)}
\author{Project Team}
\date{}

\begin{document}
\maketitle

\section*{Project Deliverable Structure}
The project is organized into three core deliverables that incrementally lead to the final output:
\begin{enumerate}[noitemsep]
  \item \textbf{Deliverable 1 (11th February):} Functional prototype + trained disease model baseline.
  \item \textbf{Deliverable 2 (28th February):} Integrated advisory platform with analytics and preliminary results.
  \item \textbf{Deliverable 3 (Final stage):} Consolidated, validated system ready for broader field use.
\end{enumerate}

The focus of this report is on Deliverable 1 and Deliverable 2.

\section*{Deliverable 1: 11th February}
\textbf{Brief description:} Improved functional prototype with disease detection, weather advisory, and core user flow operational.

\textbf{1) Objective and research alignment}
\begin{itemize}[noitemsep]
  \item Achieves objectives on image-based disease diagnosis and weather-informed guidance.
  \item Establishes the technical base for integrated advisory in later stages.
\end{itemize}

\textbf{2) Key components/features and contribution}
\begin{itemize}[noitemsep]
  \item Disease detection pipeline: image upload, preprocessing, ONNX inference, treatment mapping.
  \item Weather module: real API integration and advisory generation.
  \item Authenticated frontend workflow: login, dashboard, disease page, weather page.
  \item Contribution to final project: validates end-to-end viability of AI-assisted farmer support.
\end{itemize}

\textbf{3) Technical work involved}
\begin{itemize}[noitemsep]
  \item Prepare disease dataset (folder-per-class), create label map, split train/val/test.
  \item Train MobileNetV2-based classifier, evaluate, and export ONNX model.
  \item Integrate backend APIs with frontend screens and perform functional testing.
  \item Implement input validation to reduce unreliable detections.
\end{itemize}

\textbf{4) Scope and limitations}
\begin{itemize}[noitemsep]
  \item Scope: reliable prototype behavior on supported crops/classes.
  \item Limitations: constrained disease classes from available dataset; no field-scale pilot yet.
\end{itemize}

\textbf{5) Technology Readiness Level (TRL):} \textbf{TRL 5} (technology validated in a relevant development environment).

\section*{Deliverable 2: 28th February}
\textbf{Brief description:} Integrated system with analytical capabilities (market trends, RAG advisory), stronger robustness controls, and preliminary performance evidence.

\textbf{1) Objective and research alignment}
\begin{itemize}[noitemsep]
  \item Achieves objectives on integrated decision support: disease + weather + market + advisory.
  \item Provides preliminary quantitative/functional evidence for feasibility.
\end{itemize}

\textbf{2) Key components/features and contribution}
\begin{itemize}[noitemsep]
  \item Integrated advisory flow combining modules (disease, weather, seasonal, RAG knowledge support).
  \item Market module with realistic regional reference data and trend/prediction endpoints.
  \item Robust inference controls (confidence/margin/entropy, crop-gated checks, cautious fallback).
  \item Contribution to final project: transitions prototype to a coherent analytical advisory system.
\end{itemize}

\textbf{3) Technical work involved}
\begin{itemize}[noitemsep]
  \item Refine model inference logic for real-world robustness and reduced false guidance.
  \item Standardize weather outputs (including wind speed and units) with reliable fallback.
  \item Integrate and test market analytics endpoints and frontend visualization.
  \item Consolidate preprocessing and implementation documentation for reproducibility.
\end{itemize}

\textbf{4) Scope and limitations}
\begin{itemize}[noitemsep]
  \item Scope: integrated platform with preliminary analytical outputs and improved reliability.
  \item Limitations: still pre-pilot; model generalization depends on diversity of training images and requires continued calibration with real user data.
\end{itemize}

\textbf{5) Technology Readiness Level (TRL):} \textbf{TRL 6} (technology demonstrated as an integrated system in a relevant environment).

\section*{How Incremental Deliverables Lead to Final Outcome}
\begin{itemize}[noitemsep]
  \item \textbf{Deliverable 1} proves core feasibility (functional AI prototype).
  \item \textbf{Deliverable 2} adds integration depth and analytical readiness (system-level intelligence).
  \item \textbf{Deliverable 3} builds on both to produce a validated, deployment-oriented outcome.
\end{itemize}

\begin{table}[h]
\centering
\small
\begin{tabular}{@{}lll@{}}
\toprule
\textbf{Item} & \textbf{11th February} & \textbf{28th February} \\
\midrule
Deliverable type & Functional prototype + trained model & Integrated analytical system + preliminary results \\
Primary objective & Disease + weather feasibility & Full advisory integration and robustness \\
TRL & 5 & 6 \\
\bottomrule
\end{tabular}
\end{table}

\end{document}

